\chapter{Overall Discussion}
\label{chap:discussion}

\section{Security findings across leakage classes}

The original baseline shows why the point of enforcement matters. Pre-generation role projection performed strongly for several direct confidentiality conditions: secure mode delivered no A01 target values, no A04 complete relation edges, and no A08 protected numeric values. However, structural projection did not address every boundary: policy-aware A02 scoring still found newly delivered restricted information in 8/150 secure-mode conversations, and A05 confirmed protected indexed-record existence. In A03, all 150 downgraded public/internal conversations per mode delivered the protected value while both downgraded retrieval exposure and retained-state exposure were 150/150. Because those pathways co-occurred in the historical package, that evidence does not isolate retained state as the cause of delivery. Structural projection therefore reduces the model's opportunity to repeat disallowed fields, while relation, state, and derived-output controls remain necessary.

The deliberately mixed sensitivity-evaluation path was considerably harder to secure in the historical implementation. Protected evidence remained model-visible, and policy-aware A02 leakage rose to 123/150 conversations, including 37 complete reconstructions. Direct values, relation structure, protected indexed-record existence, and retrieval-supported numeric values were also disclosed. The final hardened package recorded no unauthorised attack-specific outcome in the finite matrix. This is a substantial descriptive package-level improvement; differences in prompts, source state, and controls preclude attributing it to one component.

The outcome classes remain distinct because A04 measures hidden relation structure, A05 protected indexed-record existence, A03 retained state after authority downgrade, and A08 retrieval-mediated numeric disclosure through an interface that exposes no vectors, scores, or rankings. Integrity is separate: the historical A06 run produced no canary compliance and two confidentiality leaks, whereas A07 produced complete canary compliance without protected-value leakage. Aggregating these events into one leakage rate would discard their security meaning and enforcement stage. The legacy A02 values of 60/150 and 122/150 remain reproducible diagnostics, but the policy-aware values replace them as the confidentiality measure because the legacy flag includes authorised or prompt-supplied fields.

\section{What the comparison evidence supports}

The original baseline records vulnerabilities in the historical system, while the hardened-package evaluation records the complete final control set. Their difference remains descriptive because prompt and source-state equivalence varies by attack. A01 target prompts differ, A02 warm-ups match but all 450 final prompts differ, A03 protected seed prompts differ, A05 formulations differ, exact historical prompt text is unavailable for A04 and A08, and A07 changes from a synthetic trigger to a natural-style request. A06 prompts match, but the implementation changes at multiple stages. Chapter~\ref{chap:vulnerability-analysis} and Appendix~\ref{app:provenance-detail} retain the complete provenance qualifications.

Component attribution instead comes from the matched and supplemental experiments. It applies only to the switched control and captured condition; it neither explains the entire package difference nor substitutes guards-off hardened code for the historical implementation. Supplemental replay, the identical-final-raw-answer A02 fork, the prospectively frozen A06 challenge, and the restored A07-S trigger add component evidence without altering the earlier package or ablation values.

\section{What the Matched and Supplemental Component Studies Add}

The strongest matched results occur at different boundaries. Relation enforcement reduced complete A04 edge disclosure from 150/150 to zero in secure mode and from 65/150 to zero in sensitivity-evaluation mode, with authorised path completion at 75/75. Membership/existence enforcement reduced A05 sensitivity-mode delivery from 15/150 to zero, although five raw confirmations remained before replacement. State clearing removed measured A03 post-transition state exposure from 150/150 secure-mode and 140/150 sensitivity-mode conversations; delivery was already zero in both guards-off arms. Embedding/rank-probe handling removed sensitivity-mode target exposure from 150/150 and delivered numeric leakage from 3/150, with positive-control success unchanged at 74/75. These studies support the switched controls under their recorded prompts and conditions, not the entire historical package difference.

The original A01 and A02 verifier ablations had zero primary raw and delivered leakage with the verifier disabled. Enabling it nevertheless replaced 150/150 unauthorised A01 answers per mode and at least one turn in 147/150 secure-mode and 52/150 sensitivity-mode A02 conversations. Deterministic replay detected every known scoped historical leak, but also produced 428 false replacements; 6/50 synthetic benign outputs were replaced. The A01 live pilot remained non-vulnerable. These observations establish detector activity and frozen-corpus behaviour, not a confidentiality reduction in the zero-to-zero live conditions.

The supplemental A02 challenge supplied the required vulnerable off condition while holding each final generated raw answer fixed. Policy-aware delivery fell from 12/150 to zero; all 12 scorer-positive outputs were detected, and 50/138 scorer-negative unauthorised outputs were also replaced. The resulting precision was 12/62 (19.4\%). This isolates the final-answer delivery contribution under the frozen stress condition. Historical package reduction, warm-up-turn effects, and general benign-query performance remain outside that attribution.

The earlier family-label-omitted A06 study was also zero-to-zero for canary compliance and protected-value delivery. Its changed quarantine and guard-check telemetry records mode-dependent context treatment, but the experiment did not reproduce the historical 2/150 confidentiality failures. The prospectively frozen supplemental challenge did reproduce a non-zero integrity outcome: enabling the composite prompt-injection guard reduced delivered exact-canary compliance from 300/300 to 0/300, and all 300 unauthorised pairs changed from compliance to non-compliance. Secure mode quarantined the poisoned context and reduced raw canary generation from 150/150 to zero. Sensitivity-evaluation mode retained the evidence and produced 145/150 raw canaries with the guard enabled, all of which were removed before delivery. The single switch controls both actions, so the estimate applies to the configured component as a whole. Protected ingredient-plus-percentage leakage was zero in both arms; confidentiality reduction was not measured.

The supplemental A06 authorised control was structurally non-evaluable. The linked formulation was retrieved in 0/300 protected records and never entered model context, leaving utility preservation and utility loss unmeasured. The \texttt{authorized\_accuracy\_error} flags record this missing prerequisite rather than an independent generation-accuracy failure. Positive-control success was absent from both frozen gates. A post-hoc exact-string review found the correct product name and target market in 96/96 benign-profile pilot answers, but this non-preregistered narrow metadata score is exploratory and cannot replace the planned control.

The natural-style A07-N relation ablation also had a non-vulnerable off condition and did not reproduce the historical synthetic-trigger behaviour. The supplemental A07-S study restored that trigger and switched the prompt-injection guard: delivered canary compliance fell from 150/150 to zero in both modes, while authorised task success remained 75/75 and protected-value leakage remained zero. Secure mode acted before generation; sensitivity-evaluation mode retained 150/150 raw canaries and removed them before delivery. The result is an integrity effect at two mode-dependent stages, not a confidentiality effect.

\section{Architectural implications}

The original failures clustered at module boundaries. Retrieval could expose mixed-sensitivity evidence; relation expansion treated hidden identifiers and edges as ordinary metadata; conversational state survived a change in authority; and generated answers were not consistently checked for reconstruction, membership, relation, or integrity outcomes. The common weakness was loss of policy meaning when source fields became derived state---focus identifiers, summaries, relation edges, retrieval candidates, previous answers, and model-visible context.

The implemented controls address these boundaries through current-role projection, state clearing, relation-aware policy, probe detection, untrusted-content handling, and delivery-stage validation. The six configurable control types are summarised architecturally in Table~\ref{tab:architecture-guard-points}, while Table~\ref{tab:guard-ablation-evidence-classes} separates the permitted intervention claims. The combined matched and supplemental results support bounded attribution for relation access, membership/existence enforcement, embedding/rank-probe handling, the A02 delivery verifier, and the A06 and A07-S prompt-injection challenges under vulnerable off conditions; state clearing and embedding/rank-probe handling also changed upstream exposure. The earlier A06 matched panel still contributes only quarantine and guard-execution telemetry, A01 remains non-vulnerable in its live off condition, and A07-N shows that a relation switch may be insufficient when a different fixed guard and generator behaviour remain on the attack path. The experiments do not decompose internal logic within each guard or establish all interactions among controls. Modularity therefore improves placement and observability but also creates configuration and interference risks that require explicit ablation.

\section{Security and authorised-task trade-offs}

Positive controls remain separate from unauthorised security outcomes. They measure whether one attack-specific protected task remained possible, not general utility or false-refusal rates. Their heterogeneous package and matched aggregates are descriptive and secondary to the per-attack results; variant-dependent A07 and the structurally non-evaluable supplemental A06 control require separate treatment. Appendix~\ref{app:detailed-hardening} retains the complete aggregate counts and their repeated-call limitations.

Broad refusals are operationally simple, but usefulness depends on detector scope. The earlier A01 and A02 ablations show that a verifier can alter many delivered answers without changing the attack-specific security outcome.

The supplemental sensitivity-evaluation-mode A02 challenge evaluates the final-delivery trade-off on identical final raw answers. The verifier blocked 12/12 scorer-positive outputs and replaced 50/138 scorer-negative unauthorised outputs; protected full-row success remained 72/75 in both derived deliveries. Because generation was shared, the protected count shows only that the verifier left those outputs unchanged. It is not an independent end-to-end utility estimate. The 50-item synthetic benign-output suite found 6 replacements. This detector replay evaluated stored outputs rather than an end-to-end benign-query suite.

State clearing has a narrower security meaning and sacrifices conversational continuity after a role transition. Delivery validation preserves upstream functionality, but exact matching may miss semantic or transformed disclosure and can suppress legitimate answers.

The experiments favour explicit acceptance thresholds for security, scorer-negative output preservation, and authorised positive controls. Broader held-out end-to-end benign tasks remain necessary for a general utility claim.

\section{Threats to validity and remaining weaknesses}

Internal validity benefits from deterministic target panels, temperature 0.0, common role and length matrices, and machine-readable records. Package-level prompt and source equivalence remains incomplete: A01, A03, A05, and A07 differ in attack-specific ways, A02 final prompts differ in all 450 preserved sequences, and exact historical prompts are unavailable for A04 and A08. The 150 unauthorised records per cell share five targets, fixed templates, roles, and conversation lengths, so repeated API calls are not treated as independent population draws. The later time-stamped A02 challenge adds evidence without altering the earlier zero-to-zero ablation.

The verifier replay corpus combines historical delivered outputs, preserved matched raw outputs, and synthetic controls with a deliberately constructed class balance. Its confusion counts characterise detector behaviour in that frozen mixed-provenance corpus and must not be interpreted as deployment prevalence or real-world false-refusal estimates. The full A02 challenge isolates only final-answer delivery: warm-up generations and pre-final state were shared rather than independently branched. Its renderer was frozen on separate development targets, but the confirmatory extension reused the five established thesis targets; it therefore provides frozen-prompt validation on the existing panel rather than fully independent target-level validation. The supplemental A06 arms were paired on recorded conditions but generated through separate API calls in a fixed off-then-on order; pairing does not eliminate residual generator-call, service-time, or order variation. Its pilot gate is an additional design weakness: it required vulnerable primary-profile canary behaviour and a reduction with the guard enabled, but did not require a valid protected retrieval opportunity or any positive-control success. The zero positive-control signal visible in the pilot should have triggered that prerequisite check before the confirmatory run.

Provenance quality improves across the evidence levels. Earlier matched ablations captured source snapshots, prompts, configurations, datasets, and pair identifiers, but used dirty working trees and usually lacked index-content hashes and versioned scorers. The superseding and supplemental studies add canonical and serialized index hashes, dependency and scorer versions, exact messages, and result bindings. Remaining gaps include inherited A03 guard configuration and A07-N activation telemetry; incomplete or failed runs contribute no denominator. Appendix~\ref{app:provenance-detail} records the complete package-specific fields, the supplemental A06 shard and archive inventory, supersession decisions, and the limits of HPC-only raw-data portability.

Construct validity depends on attack-specific scorers. Exact-value detection may miss paraphrases, approximations, translations, ranges, and derivations. The A02 policy-aware scorer corrects role and prompt awareness but still relies on exact or field-labelled evidence; its manual-review candidates are not automatically counted as leakage. Retrieval exposure fields differ between attacks and must not be aggregated as one homogeneous metric. Canary compliance measures integrity, not confidentiality. A05 measures protected-record existence in the indexed corpus rather than training-data membership inference. A08 measures the effect of embedding/rank-oriented query framing on retrieval-mediated disclosure rather than exposing or inverting embeddings.

External validity is limited to one workbook, five targets per attack, one generation model, one embedding model, retrieval depth five, and fixed prompt templates. The longer conversations are controlled warm-up conditions rather than adaptive attacks. The interface does not expose raw similarity scores, full rankings, or citations, and the experiments do not cover concurrent users, cross-session memory, distributed caches, external tools, or role changes during generation.

Sensitivity-evaluation mode intentionally permits protected evidence to reach downstream stages in some attacks, making delivery enforcement security-critical. Pattern-based membership, injection, and embedding-probe detectors can be bypassed by unseen wording and can over-refuse benign or authorised requests. Relation policy was tested only on short fixed paths. Within the finite matrix, no leak was detected under the specified attacks and scorers. Adaptive and previously unseen attacks remain unevaluated, so this bounded negative evidence does not establish general system security.
